\Needspace{16\baselineskip}
\Article{13}

Die Faktoren auf der rechten Seite der Gleichung $[5]$ lassen sich auch folgendermaßen schreiben:

\[
\begin{aligned}
r-r^{-1}&=-(r^{n-1}-r^{-n+1})\\
r^3-r^{-3}&=-(r^{n-3}-r^{-n+3})\\
r^5-r^{-5}&=-(r^{n-5}-r^{-n+5}),\;\text{usw.}
\end{aligned}
\]

bis hin zu

\[
\begin{aligned}
r^{n-2}-r^{-n+2}=-(r^2-r^{-2})
\end{aligned}
\]

Damit nimmt die Gleichung diese Form an:

\SourceLeaf{35}

\[
\begin{aligned}
W&=(-1)^{\frac12(n-1)}(r^2-r^{-2})(r^4-r^{-4})(r^6-r^{-6})\\&\qquad\dotsfour(r^{n-1}-r^{-n+1})
\end{aligned}
\]

Multiplizieren wir diese Gleichung mit $[5]$ in ihrer ursprünglichen Form, so ergibt sich

\[
\begin{aligned}
W^2&=(-1)^{\frac12(n-1)}(r-r^{-1})(r^2-r^{-2})(r^3-r^{-3})\\&\qquad\dotsfour(r^{n-1}-r^{-n+1})
\end{aligned}
\]

wobei $(-1)^{\frac12(n-1)}$ entweder $=+1$ oder $=-1$ ist, je nachdem, ob $n$ die Form $4\mu+1$ oder die Form $4\mu+3$ hat.\quad Daraus folgt

\[
\begin{aligned}
W^2&=\pm r^{\frac12n(n-1)}(1-r^{-2})(1-r^{-4})(1-r^{-6})\\&\qquad\dotsfour(1-r^{-2(n-1)})
\end{aligned}
\]

Nun sieht man ohne Mühe, dass $r^{-2}$, $r^{-4}$, $r^{-6}\dotsfour r^{-2n+2}$ sämtliche Wurzeln der Gleichung $x^n-1=0$ mit Ausnahme der Wurzel $x=1$ darstellen. Daher muss die folgende Gleichung identisch für eine Unbestimmte gelten:

\[
\begin{aligned}
&(x-r^{-2})(x-r^{-4})(x-r^{-6})\dotsfour(x-r^{-2n+2})\\&\qquad=x^{n-1}+x^{n-2}+x^{n-3}+\text{usw.}+x+1
\end{aligned}
\]

Setzen wir also $x=1$, so erhalten wir

\[
\begin{aligned}
(1-r^{-2})(1-r^{-4})(1-r^{-6})\dotsfour(1-r^{-2n+2})=n
\end{aligned}
\]

und da offenbar $r^{\frac12n(n-1)}=1$ ist, geht unsere Gleichung in die folgende über:

\[
\begin{aligned}
W^2=\pm n\  .\  .\  .\  .\  .\  .\  .\  .\  .\  .\ [6]
\end{aligned}
\]

In dem Fall, dass $n$ die Form $4\mu+1$ hat, ist somit

\[
\begin{aligned}
W=\pm\surd n,\quad\text{und somit}\quad T=\pm\surd n,\quad U=0
\end{aligned}
\]

Im anderen Fall dagegen, in dem $n$ die Form $4\mu+3$ hat, ist

\[
\begin{aligned}
W=\pm i\surd n,\quad\text{und somit}\quad T=0,\quad U=\pm\surd n
\end{aligned}
\]

\Needspace{13\baselineskip}
\Article{14}

Die Methode des vorigen Artikels liefert nur den absoluten Wert der Summen $T$, $U$ und lässt offen, ob für $T$ im ersten Fall und für $U$ im zweiten Fall $=+\surd n$ oder $=-\surd n$ zu setzen ist.\quad Dies lässt sich aber, wenigstens für den Fall $k=1$, anhand der Gleichung $[5]$ folgendermaßen entscheiden.\quad Da für $k=1$

\SourceLeaf{36}

\[
\begin{aligned}
r-r^{-1}&=2i\sin\omega\\
r^3-r^{-3}&=2i\sin3\omega\\
r^5-r^{-5}&=2i\sin5\omega\;\text{usw.}
\end{aligned}
\]

ist, geht jene Gleichung über in

\[
\begin{aligned}
W=(2i)^{\frac12(n-1)}\sin\omega\sin3\omega\sin5\omega\ldots\sin(n-2)\omega
\end{aligned}
\]

In dem Fall, dass $n$ die Form $4\mu+1$ hat, finden sich nun in der Folge der ungeraden Zahlen

\[
\begin{aligned}
1,\;3,\;5,\;7\dotsfour\tfrac12(n-3),\quad\tfrac12(n+1)\dotsfour(n-2)
\end{aligned}
\]

$\frac14(n-1)$ Zahlen, die kleiner als $\frac12n$ sind und denen offenbar positive Sinuswerte entsprechen. Die übrigen $\frac14(n-1)$ dagegen sind größer als $\frac12n$, und ihnen entsprechen negative Sinuswerte. Das Produkt sämtlicher Sinuswerte ist daher gleich einer positiven Größe, multipliziert mit $(-1)^{\frac14(n-1)}$. Somit ist $W$ gleich einer positiven reellen Größe, multipliziert mit $i^{n-1}$ oder mit $1$, da $i^4=1$ ist und $n-1$ durch $4$ teilbar ist: Die Größe $W$ ist also reell und positiv. Daher gilt notwendig

\[
\begin{aligned}
W=+\surd n,\quad T=+\surd n
\end{aligned}
\]

Im anderen Fall, in dem $n$ die Form $4\mu+3$ hat, sind in der Folge der ungeraden Zahlen

\[
\begin{aligned}
1,\;3,\;5,\;7\dotsfour\tfrac12(n-1),\quad\tfrac12(n+3)\dotsfour(n-2)
\end{aligned}
\]

\Needspace{12\baselineskip}
die ersten $\frac14(n+1)$ kleiner als $\frac12n$, die übrigen $\frac14(n-3)$ dagegen größer.\quad Unter den Sinuswerten der Bögen $\omega$, $3\omega$, $5\omega\dotsfour(n-2)\omega$ sind somit $\frac14(n-3)$ negativ, und daher ist $W$ das Produkt aus $i^{\frac12(n-1)}$, einer positiven reellen Größe und $(-1)^{\frac14(n-3)}$. Der dritte Faktor ist $=i^{\frac12(n-3)}$ und ergibt mit dem ersten multipliziert $i^{n-2}=i$, da $i^{n-3}=1$ ist.\quad Somit gilt notwendig

\[
\begin{aligned}
W=+i\surd n,\quad\text{und}\quad U=+\surd n
\end{aligned}
\]

\Needspace{15\baselineskip}
\Article{15}

Wir werden nun zeigen, wie sich dieselben Ergebnisse aus der in Art. $9$ betrachteten Reihe ableiten lassen.\quad Ersetzen wir in der Gleichung $[4]$ den Ausdruck $x^{\frac12}$ durch $-y^{-1}$, so erhalten wir

\SourceLeaf{37}

\[
\begin{aligned}
&1-y^{-1}\frac{1-y^{-2m}}{1-y^{-2}}+y^{-2}\frac{(1-y^{-2m})(1-y^{-2m+2})}{(1-y^{-2})(1-y^{-4})}\\&\qquad-y^{-3}\frac{(1-y^{-2m})(1-y^{-2m+2})(1-y^{-2m+4})}{(1-y^{-2})(1-y^{-4})(1-y^{-6})}+\text{usw.}
\end{aligned}
\]

bis zum Glied an der Stelle $m+1$

\[
\begin{aligned}
&=(1-y^{-1})(1+y^{-2})(1-y^{-3})(1+y^{-4})\\&\qquad\ldots(1\pm y^{-m})\  .\  .\  .\  .\ [7]
\end{aligned}
\]

Wählen wir hier für $y$ eine eigentliche Wurzel der Gleichung $y^n-1=0$, etwa $r$, und setzen zugleich $m=n-1$, so ist

\[
\begin{aligned}
\frac{1-y^{-2m}}{1-y^{-2}}&=\frac{1-r^2}{1-r^{-2}}=-r^2\\
\frac{1-y^{-2m+2}}{1-y^{-4}}&=\frac{1-r^4}{1-r^{-4}}=-r^4\\
\frac{1-y^{-2m+4}}{1-y^{-6}}&=\frac{1-r^6}{1-r^{-6}}=-r^6
\end{aligned}
\]

bis hin zu

\[
\begin{aligned}
\frac{1-y^{-2}}{1-y^{-2m}}=\frac{1-r^{2n-2}}{1-r^{-2n+2}}=-r^{2n-2}
\end{aligned}
\]

Dabei ist zu beachten, dass keiner der Nenner $1-r^{-2}$, $1-r^{-4}$ usw. $=0$ wird.\quad Daher nimmt die Gleichung $[7]$ diese Form an:

\[
\begin{aligned}
&1+r+r^4+r^9+\text{usw.}+r^{(n-1)^2}\\&\qquad=(1-r^{-1})(1+r^{-2})(1-r^{-3})\dotsfour(1+r^{-n+1})
\end{aligned}
\]

\Needspace{10\baselineskip}
Multiplizieren wir auf der rechten Seite dieser Gleichung das erste Glied mit dem letzten, das zweite mit dem vorletzten usw., so erhalten wir

\[
\begin{aligned}
(1-r^{-1})(1+r^{-n+1})&=r-r^{-1}\\
(1+r^{-2})(1-r^{-n+2})&=r^{n-2}-r^{-n+2}\\
(1-r^{-3})(1+r^{-n+3})&=r^3-r^{-3}\\
(1+r^{-4})(1-r^{-n+4})&=r^{n-4}-r^{-n+4}\;\text{usw.}
\end{aligned}
\]

Man sieht leicht, dass diese Teilprodukte das folgende Produkt ergeben:

\[
\begin{aligned}
&(r-r^{-1})(r^3-r^{-3})(r^5-r^{-5})\\&\qquad\dotsfour(r^{n-4}-r^{-n+4})(r^{n-2}-r^{-n+2})
\end{aligned}
\]

Dieses ist somit

\[
\begin{aligned}
=1+r+r^4+r^9+\text{usw.}+r^{(n-1)^2}=W
\end{aligned}
\]

Diese Gleichung ist identisch mit der Gleichung $[5]$, die in Art. $12$ aus der ersten Reihe abgeleitet wurde; die übrigen Schlussfolgerungen ergeben sich dann auf dieselbe Weise wie in den Art. $13$ und $14$.

\SourceLeaf{38}

\Needspace{16\baselineskip}
\Article{16}

Wir gehen zum anderen Fall über, in dem $n$ eine gerade Zahl ist.\quad Zunächst habe $n$ die Form $4\mu+2$, sei also das Doppelte einer ungeraden Zahl. Offenbar liefern die Zahlen $\frac14nn$, $(\frac12n+1)^2-1$, $(\frac12n+2)^2-4$ usw., oder allgemein $(\frac12n+\lambda)^2-\lambda\lambda$, bei Division durch $\frac12n$ ungerade Quotienten und sind daher modulo $n$ zu $\frac12n$ kongruent.\quad Daraus folgt, dass für eine eigentliche Wurzel $r$ der Gleichung $x^n-1=0$, für die also $r^{\frac12n}=-1$ ist, gilt:

\[
\begin{aligned}
r^{(\frac12n)^2}&=-1\\
r^{(\frac12n+1)^2}&=-r\\
r^{(\frac12n+2)^2}&=-r^4\\
r^{(\frac12n+2)^2}&=-r^9\;\text{usw.}
\end{aligned}
\]

Somit hebt in der Reihe

\[
\begin{aligned}
1+r+r^4+r^9+\text{usw.}+r^{(n-1)^2}
\end{aligned}
\]

das Glied $r^{(\frac12n)^2}$ das erste auf, das nächste das zweite usw., sodass sich ergibt:

\[
\begin{aligned}
W=0,\quad T=0,\quad U=0
\end{aligned}
\]

\Needspace{9\baselineskip}
\Article{17}

Es bleibt der Fall, in dem $n$ die Form $4\mu$ hat, also durch vier teilbar ist.\quad Hier ist allgemein $(\frac12n+\lambda)^2-\lambda\lambda$ durch $n$ teilbar, und somit gilt

\[
\begin{aligned}
r^{(\frac12n+\lambda)^2}=r^{\lambda\lambda}
\end{aligned}
\]

In der Reihe

\[
\begin{aligned}
1+r+r^4+r^9+\text{usw.}+r^{(n-1)^2}
\end{aligned}
\]

ist daher das Glied $r^{(\frac12n)^2}$ dem ersten gleich, das nächste dem zweiten usw., sodass sich ergibt:

\[
\begin{aligned}
W=2(1+r+r^4+r^9+\text{usw.}+r^{(\frac12n-1)^2})
\end{aligned}
\]

Nehmen wir nun an, dass in der Gleichung $[7]$ aus Art. 15 $m=\frac12n-1$ gesetzt und für $y$ eine eigentliche Wurzel der Gleichung $y^n-1=0$, etwa $r$, gewählt wird.\quad Dann erhält die Gleichung, ebenso wie in Art. 15, die folgende Form:

\[
\begin{aligned}
&1+r+r^4+\text{usw.}+r^{(\frac12n-1)^2}\\&\qquad=(1-r^{-1})(1+r^{-2})(1-r^{-3})\dotsfour(1-r^{-\frac12n+1})
\end{aligned}
\]

\SourceLeaf{39}

oder

\[
\begin{aligned}
W&=2(1-r^{-1})(1+r^{-2})(1-r^{-3})(1+r^{-4})\\&\qquad\ldots(1-r^{-\frac12n+1})\quad\dotstwo\quad[8]
\end{aligned}
\]

\Needspace{10\baselineskip}
Da ferner $r^{\frac12n}=-1$ und folglich

\[
\begin{aligned}
1+r^{-2}&=-r^{\frac12n-2}(1-r^{-\frac12n+2})\\
1+r^{-4}&=-r^{\frac12n-4}(1-r^{-\frac12n+4})\\
1+r^{-6}&=-r^{\frac12n-6}(1-r^{-\frac12n+6})\;\text{usw.}
\end{aligned}
\]

ist und das Produkt der Faktoren $-r^{\frac12n-2}$, $-r^{\frac12n-4}$, $-r^{\frac12n-6}$ usw. bis zu $-r^2$ gleich $=(-1)^{\frac14n-1}\,r^{\frac1{16}nn-\frac14n}$ ist, lässt sich die vorige Gleichung auch folgendermaßen schreiben:

\[
\begin{aligned}
W&=2(-1)^{\frac14n-1}\,r^{\frac1{16}nn-\frac14n}(1-r^{-1})(1-r^{-2})\\&\qquad(1-r^{-3})(1-r^{-4})\dotsfour(1-r^{-\frac12n+1})
\end{aligned}
\]

\Needspace{10\baselineskip}
Da

\[
\begin{aligned}
1-r^{-1}&=-r^{-1}(1-r^{-n+1})\\
1-r^{-2}&=-r^{-2}(1-r^{-n+2})\\
1-r^{-3}&=-r^{-3}(1-r^{-n+3})\;\text{usw.}
\end{aligned}
\]

ist, erhalten wir

\[
\begin{aligned}
&(1-r^{-1})(1-r^{-2})(1-r^{-3})\ldots(1-r^{-\frac12n+1})\\
&\quad=(-1)^{\frac12n-1}\,r^{-\frac18nn+\frac14n}(1-r^{-\frac12n-1})\\&\qquad(1-r^{-\frac12n-2})(1-r^{-\frac12n-3})\ldots(1-r^{-n+1})
\end{aligned}
\]

und damit

\[
\begin{aligned}
W&=2(-1)^{\frac34n-2}\,r^{-\frac1{16}nn}(1-r^{-\frac12n-1})\\&\qquad(1-r^{-\frac12n-2})(1-r^{-\frac12n-3})\dotsfour(1-r^{-n+1})
\end{aligned}
\]

Multiplizieren wir diesen Wert von $W$ mit dem zuvor gefundenen und fügen auf beiden Seiten den Faktor $1-r^{-\frac12n}$ hinzu, so ergibt sich

\[
\begin{aligned}
(1-r^{-\frac12n})W^2&=4(-1)^{n-3}r^{-\frac14n}(1-r^{-1}\mathord{.})\\&\qquad(1-r^{-2})(1-r^{-3})\dotsfour(1-r^{-n+1})
\end{aligned}
\]

Es ist aber

\[
\begin{aligned}
1-r^{-\frac12n}&=2\\
(-1)^{n-3}&=-1\\
r^{-\frac14n}&=-r^{\frac14n}
\end{aligned}
\]

\[
\begin{aligned}
(1-r^{-1})(1-r^{-2})(1-r^{-3})\dotsfour(1-r^{-n+1})=n
\end{aligned}
\]

\Needspace{10\baselineskip}
Daraus folgt schließlich

\SourceLeaf{40}

\[
\begin{aligned}
W^2=2r^{\frac14n}n\  .\  .\  .\  .\  .\  .\  .\  .\  .\ [9]
\end{aligned}
\]

\Needspace{14\baselineskip}
Nun sieht man leicht, dass $r^{\frac14n}$ entweder $=+i$ oder $=-i$ ist, je nachdem, ob $k$ die Form $4\mu+1$ oder die Form $4\mu+3$ hat.\quad Da ferner

\[
\begin{aligned}
2i=(1+i)^2,\quad -2i=(1-i)^2
\end{aligned}
\]

ist, ergibt sich in dem Fall, dass $k$ die Form $4\mu+1$ hat,

\[
\begin{aligned}
W=\pm(1+i)\surd n,\quad\text{und somit}\quad T=U=\pm\surd n
\end{aligned}
\]

im anderen Fall dagegen, in dem $k$ die Form $4\mu+3$ hat,

\[
\begin{aligned}
W=\pm(1-i)\surd n,\quad\text{und somit}\quad T=-U=\pm\surd n
\end{aligned}
\]
